%% =====================================================================
%%  Unit 2 — Where are you from?  (Dire d'où l'on vient)
%%  Pistes : 2.01–2.04 dialogues · 2.05–2.06 vocabulaire · 2.07–2.08 grammaire
%%           2.09 prononciation · 2.10–2.13 exercices (textes en annexe)
%% =====================================================================

\newcommand{\otIllusUnII}{\otIllusDuo{France/104/86/15/0.10,USA/190/92/20/0.09,UK/122/14/13/0.08,Japan/150/8/11/0.08,
  Brazil/196/22/12/0.07,Italy/100/50/12/0.07}{jack}{Where are you from?}{lea}{I'm from France!}}

\unite[Where are\\you from?]{2}{Where are you from?}{Dire d'où l'on vient}{otIllusUnII}

\situation{Épisode 2}{Mardi, 13 heures : c'est la pause déjeuner. Jack et Emma emmènent Léa dans
leur café préféré, à deux pas du studio. Autour d'un sandwich, les questions fusent : d'où
viens-tu ? Quelles langues parles-tu ? Le soir, à Brooklyn, Léa fait plus ample connaissance
avec sa propriétaire, Mrs.\ Lopez.}

\begin{objectifs}
\item dire de quel pays et de quelle ville vous venez ;
\item donner votre nationalité et dire quelles langues vous parlez ;
\item dire ce que vous n'êtes pas, avec la forme négative de \en{be} ;
\item poser une question avec \en{be} et y répondre : \en{Yes, I am. / No, I'm not.} ;
\item placer l'accent sur la bonne syllabe : \en{Japan}, \en{Japanese}.
\end{objectifs}

\auprogramme
  {Dire d'où l'on vient, sa nationalité, les langues que l'on parle}
  {\en{be} : la négation\par\vspace{0.6mm}Questions et réponses courtes}
  {Pays, nationalités, langues : 29 expressions}
  {L'accent de mot}
  {Le monde anglophone\par\vspace{0.6mm}Royaume-Uni ou Angleterre ?}

\semaine{casque/otNavy/Dialogues/Écouter et lire les deux scènes,
  bulle/otGreenDark/Vocabulaire/Apprendre les 29 expressions,
  blocs/otViolet/Grammaire/Négation et questions avec be,
  sons/otOrangeDark/Sons et culture/L'accent de mot et le monde anglophone,
  crayon/otNavy/Exercices/12 exercices corrigés,
  coche/otGreenDark/Bilan/L'essentiel et le test express}

\clearpage
%% ---------------------------------------------------------------------
\rubrique{ecouter}{À la pause déjeuner}

\begin{dialogue}{1}{Where are you from?}{13 heures, dans un café près du studio. Jack, Emma
  et Léa déjeunent ensemble.}{2.01}{2.02}
\dit{jack}{So, Léa, \hl{where are you from}?\note{1}}{Alors, Léa, tu viens d'où ?}
\dit{lea}{\hl{I'm from France}.}{Je viens de France.}
\dit{jack}{\hl{Are you from Paris}?}{Tu es de Paris ?}
\dit{lea}{\hl{No, I'm not}. I'm from Lyon.\note{2}}{Non. Je suis de Lyon.}
\dit{jack}{\hl{Is Lyon} a big city?\note{3}}{Lyon, c'est une grande ville ?}
\dit{lea}{\hl{Yes, it is}. It's a beautiful city!}{Oui. C'est une très belle ville !}
\dit{lea}{And you, Jack? Are you from New York?}{Et toi, Jack ? Tu es de New York ?}
\dit{jack}{\hl{Yes, I am}. I'm from Brooklyn.}{Oui. Je suis de Brooklyn.}
\dit{lea}{And Emma, you're British, \hl{right}?\note{4}}{Et toi, Emma, tu es britannique, c'est ça ?}
\dit{emma}{\hl{That's right}. I'm from London.}{Exactement. Je suis de Londres.}
\dit{emma}{\hl{I'm not American}\ldots{} but I love New York!}{Je ne suis pas américaine… mais j'adore New York !}
\end{dialogue}

\begin{comprendre}
\pointnote{1} \en{\textbf{Where are you from?}} Mot à mot : « Où es-tu de ? ». \en{from} indique
  l'origine. Pour répondre : \en{I'm from} + pays ou ville : \en{I'm from France}, \en{I'm from Lyon}.
\pointnote{2} \en{\textbf{No, I'm not.}} En anglais, on répond rarement par un simple \en{yes} ou
  \en{no} : on ajoute une \textbf{réponse courte}, qui reprend le verbe. \en{I'm not} =
  \en{I am not} (« je ne suis pas »).
\pointnote{3} \en{\textbf{Is Lyon a big city?}} Pour poser une question, on place \en{be}
  \textbf{avant} le sujet : \en{Lyon is…} devient \en{Is Lyon…?}
\pointnote{4} \en{\textbf{…, right?}} À la fin d'une phrase, \en{right?} veut dire « c'est ça ? »,
  « n'est-ce pas ? ». On confirme avec \en{That's right} (« c'est exact »).
\end{comprendre}

\begin{dialogue}{2}{I'm not Spanish}{19 heures, à Brooklyn. En rentrant, Léa entend sa propriétaire,
  Mrs.\ Lopez, parler espagnol au téléphone.}{2.03}{2.04}
\dit{lea}{Good evening, Mrs.\ Lopez. \hl{Are you Spanish}?}{Bonsoir, madame Lopez. Vous êtes espagnole ?}
\dit{lopez}{\hl{No, I'm not Spanish}. I'm American.}{Non, je ne suis pas espagnole. Je suis américaine.}
\dit{lopez}{But my parents \hl{are Mexican}.\note{5}}{Mais mes parents sont mexicains.}
\dit{lea}{Oh! So you \hl{speak Spanish}?}{Ah ! Alors vous parlez espagnol ?}
\dit{lopez}{Yes. I speak Spanish, English\ldots{} and \hl{a little} French! Bonjour!\note{6}}{Oui. Je parle espagnol, anglais… et un peu français ! Bonjour !}
\dit{lea}{Really? Your French is very good!}{Vraiment ? Votre français est très bon !}
\dit{lopez}{Thank you. And you speak English very well, Léa.}{Merci. Et vous parlez très bien anglais, Léa.}
\dit{lea}{Oh no, my English \hl{isn't} very good. I'm \hl{a beginner}!\note{7}}{Oh non, mon anglais n'est pas très bon. Je suis débutante !}
\dit{lopez}{\hl{Don't worry}. It's very good!}{Ne vous inquiétez pas. Il est très bien !}
\end{dialogue}

\begin{comprendre}
\pointnote{5} \en{\textbf{my parents are Mexican}} Les nationalités prennent toujours une
  \textbf{majuscule} en anglais : \en{Mexican}, \en{French}, \en{American}. Les langues aussi :
  \en{Spanish} veut dire à la fois « espagnol(e) » et « l'espagnol ».
\pointnote{6} \en{\textbf{I speak Spanish}} « Je parle espagnol ». \en{speak} est un verbe
  ordinaire : vous découvrirez ces verbes à l'unité 6. \en{a little} = « un peu ».
\pointnote{7} \en{\textbf{isn't}} est la forme courte de \en{is not}. \en{a beginner} = « un(e)
  débutant(e) » : en anglais, on met \en{a} devant un métier ou un statut.
\end{comprendre}

%% ---------------------------------------------------------------------
\rubrique{mots}{Les mots de l'unité}

Voici les 29 expressions de l'unité. Remarquez la majuscule des nationalités et des langues,
et l'article \en{the} devant \en{the UK} et \en{the US}, mais pas devant les autres pays.

\begin{vocabulaire}{Pays et nationalités}{2.05}
\categorie{Le monde anglophone}
\mot[England. English.]{England · English}{ˈɪŋɡlənd · ˈɪŋɡlɪʃ}{l'Angleterre · anglais(e)}
\mot[The UK. British.]{the UK · British}{ðə ˌjuːˈkeɪ · ˈbrɪtɪʃ}{le Royaume-Uni · britannique}
\mot[The US. American.]{the US · American}{ðə ˌjuːˈes · əˈmerɪkən}{les États-Unis · américain(e)}
\mot[Canada. Canadian.]{Canada · Canadian}{ˈkænədə · kəˈneɪdiən}{le Canada · canadien(ne)}
\mot[Ireland. Irish.]{Ireland · Irish}{ˈaɪrlənd · ˈaɪrɪʃ}{l'Irlande · irlandais(e)}
\mot[Australia. Australian.]{Australia · Australian}{ɑːˈstreɪljə · ɑːˈstreɪljən}{l'Australie · australien(ne)}
\categorie{Ailleurs dans le monde}
\mot[France. French.]{France · French}{fræns · frentʃ}{la France · français(e)}
\mot[Spain. Spanish.]{Spain · Spanish}{speɪn · ˈspænɪʃ}{l'Espagne · espagnol(e)}
\mot[Italy. Italian.]{Italy · Italian}{ˈɪtəli · ɪˈtæljən}{l'Italie · italien(ne)}
\mot[Germany. German.]{Germany · German}{ˈdʒɝːməni · ˈdʒɝːmən}{l'Allemagne · allemand(e)}
\mot[Mexico. Mexican.]{Mexico · Mexican}{ˈmeksɪkoʊ · ˈmeksɪkən}{le Mexique · mexicain(e)}
\mot[Brazil. Brazilian.]{Brazil · Brazilian}{brəˈzɪl · brəˈzɪliən}{le Brésil · brésilien(ne)}
\mot[China. Chinese.]{China · Chinese}{ˈtʃaɪnə · tʃaɪˈniːz}{la Chine · chinois(e)}
\mot[Japan. Japanese.]{Japan · Japanese}{dʒəˈpæn · ˌdʒæpəˈniːz}{le Japon · japonais(e)}
\end{vocabulaire}

\begin{vocabulaire}{Dire d'où l'on vient}{2.06}
\categorie{Questions et réponses}
\mot{Where are you from?}{ˌwer ɑːr juː ˈfrʌm}{D'où viens-tu ? / D'où venez-vous ?}
\mot{I'm from France.}{aɪm frəm ˈfræns}{Je viens de France.}
\mot{Are you French?}{ɑːr juː ˈfrentʃ}{Tu es français(e) ? / Vous êtes français(e) ?}
\mot{Yes, I am.}{ˌjes aɪ ˈæm}{Oui. (réponse courte)}
\mot{No, I'm not.}{ˌnoʊ aɪm ˈnɑːt}{Non. (réponse courte)}
\mot{That's right.}{ˌðæts ˈraɪt}{C'est exact.}
\mot{Really?}{ˈriːəli}{Vraiment ?}
\categorie{Langues et villes}
\mot{I speak English.}{aɪ ˌspiːk ˈɪŋɡlɪʃ}{Je parle anglais.}
\mot{a language}{ə ˈlæŋɡwɪdʒ}{une langue}
\mot{a little}{ə ˈlɪtəl}{un peu}
\mot{a beginner}{ə bɪˈɡɪnɚ}{un débutant, une débutante}
\mot{a big city}{ə ˌbɪɡ ˈsɪti}{une grande ville}
\mot{a beautiful city}{ə ˌbjuːtɪfəl ˈsɪti}{une belle ville}
\mot{my parents}{maɪ ˈperənts}{mes parents}
\mot{Don't worry.}{ˌdoʊnt ˈwɝːi}{Ne t'inquiète pas. / Ne vous inquiétez pas.}
\end{vocabulaire}

\begin{astuce}[Retenir les nationalités]
Regroupez-les par terminaison : \en{-an} / \en{-ian} (\en{American}, \en{Mexican}, \en{Italian},
\en{Canadian}, \en{Brazilian}), \en{-ish} (\en{English}, \en{British}, \en{Spanish}, \en{Irish}),
\en{-ese} (\en{Chinese}, \en{Japanese}). Il ne reste que deux exceptions : \en{French} et \en{German}.
\end{astuce}

%% ---------------------------------------------------------------------
\rubrique{grammaire}{Le verbe be : négation et questions}

\otIntertitre[otViolet]{1 · La forme négative}
Pour dire ce que l'on n'est pas, on ajoute \en{\textbf{not}} après \en{be}. À l'oral, on contracte
presque toujours : soit le sujet et le verbe (\en{you're not}), soit le verbe et \en{not}
(\en{you aren't}). Les deux sont correctes et ont le même sens.

\begin{center}\otTableau
\begin{tabular}{@{}>{\enOn}l<{\enOff}@{\hspace{8mm}}>{\enOn\bfseries}l<{\enOff}@{\hspace{8mm}}l@{}}
\otTH{Forme pleine} & \otTH{Formes contractées} & \otTH{Français}\\\midrule
I am not     & I'm not                    & je ne suis pas\\
you are not  & you're not · you aren't    & tu n'es pas, vous n'êtes pas\\
he is not    & he's not · he isn't        & il n'est pas\\
she is not   & she's not · she isn't      & elle n'est pas\\
it is not    & it's not · it isn't        & il / elle n'est pas, ce n'est pas\\\midrule
we are not   & we're not · we aren't      & nous ne sommes pas\\
you are not  & you're not · you aren't    & vous n'êtes pas\\
they are not & they're not · they aren't  & ils / elles ne sont pas\\\bottomrule
\end{tabular}\end{center}

\blocaudio{2.07}{Écoutez et répétez}{Sept phrases à la forme négative, avec les contractions
\en{isn't} \ipa{ˈɪzənt} et \en{aren't} \ipa{ɑːrnt}. Répétez pendant les pauses.}

\otIntertitre[otViolet]{2 · Les questions et les réponses courtes}
Pour poser une question, on \textbf{inverse} le sujet et le verbe \en{be} : le verbe passe en
premier. On répond par une \textbf{réponse courte} qui reprend le sujet et le verbe.

\begin{center}\otTableau
\begin{tabular}{@{}>{\enOn}l<{\enOff}@{\hspace{7mm}}>{\enOn\bfseries}l<{\enOff}@{\hspace{7mm}}>{\enOn}l<{\enOff}@{}}
\otTH{Affirmation} & \otTH{Question} & \otTH{Réponses courtes}\\\midrule
You are French.        & Are you French?        & Yes, I am. / No, I'm not.\\
He is American.        & Is he American?        & Yes, he is. / No, he isn't.\\
She is from London.    & Is she from London?    & Yes, she is. / No, she isn't.\\
It is a big city.      & Is it a big city?      & Yes, it is. / No, it isn't.\\
They are Italian.      & Are they Italian?      & Yes, they are. / No, they aren't.\\\bottomrule
\end{tabular}\end{center}

\blocaudio{2.08}{Écoutez et répétez}{Une question, puis sa réponse courte. Répétez les deux, en
faisant monter la voix à la fin de la question.}

\begin{paire}
\begin{regle}[La question avec \en{where}]
Le mot interrogatif se place en tête, puis on inverse :\par\vspace{1mm}
\exen{Where are you from?}{}
\exen{Where is she from?}{→ \en{Where's she from?}}
\exen{Where are they from?}{}
\end{regle}
\begin{regle}[\en{the} devant les pays ?]
Non, sauf pour les noms qui contiennent un nom commun :\par\vspace{1mm}
\exen{I'm from France.}{(pas d'article)}
\exen{I'm from the US.}{\en{United States}}
\exen{I'm from the UK.}{\en{United Kingdom}}
\end{regle}
\end{paire}

\begin{attention}[Les pièges des francophones]
\fauxjuste{Yes, I'm.}{Yes, I am.}
\fauxjuste{I amn't French.}{I'm not French.}
\fauxjuste{Is French Léa?}{Is Léa French?}
\fauxjuste{I'm french.}{I'm French.}
\fauxjuste{I'm from the France.}{I'm from France.}
\par\vspace{1mm}{\fontsize{10}{12.6}\selectfont Dans une réponse courte affirmative, \textbf{jamais} de
contraction : \en{Yes, I am}, \en{Yes, she is}, \en{Yes, they are}.\par}
\end{attention}

%% ---------------------------------------------------------------------
\rubrique{sons}{L'accent de mot}

En français, toutes les syllabes d'un mot ont à peu près la même force, et la voix appuie un peu
sur la dernière. En anglais, chaque mot a \textbf{une syllabe accentuée} : elle est plus
\textbf{forte}, plus \textbf{longue} et plus \textbf{haute} que les autres. Les autres syllabes
s'affaiblissent, souvent jusqu'au son \phon{ə}. Dans la phonétique, le signe \phon{ˈ} se place
juste avant la syllabe accentuée.

\begin{center}
\begin{tikzpicture}
  \foreach \s/\x/\h/\c in {ja/0/7/otMuted!60,pan/16/13/otOrangeDark}{%
    \fill[\c,rounded corners=1mm] (\x mm,0) rectangle ++(14mm,\h mm);
    \node[anchor=north,font=\otSemi\fontsize{10}{12}\selectfont] at (\x mm+7mm,-1mm) {\enb{\s}};}
  \node[anchor=north,font=\otRegular\fontsize{8.4}{10}\selectfont,text=otMuted] at (15mm,-6.5mm) {\enb{Japan} \ipa{dʒəˈpæn}};
  \foreach \s/\x/\h/\c in {ja/55/9/otMuted!60,pa/71/6/otMuted!60,nese/87/13/otOrangeDark}{%
    \fill[\c,rounded corners=1mm] (\x mm,0) rectangle ++(14mm,\h mm);
    \node[anchor=north,font=\otSemi\fontsize{10}{12}\selectfont] at (\x mm+7mm,-1mm) {\enb{\s}};}
  \node[anchor=north,font=\otRegular\fontsize{8.4}{10}\selectfont,text=otMuted] at (78mm,-6.5mm) {\enb{Japanese} \ipa{ˌdʒæpəˈniːz}};
  \node[anchor=west,text width=52mm,font=\fontsize{9.4}{11.6}\selectfont,align=left] at (110mm,7mm)
    {La hauteur des barres montre la force de chaque syllabe : l'accent \textbf{change de place} entre le pays et la nationalité.};
\end{tikzpicture}
\end{center}

\begin{paires}{2.09}{La syllabe accentuée est soulignée. Écoutez le pays, puis la nationalité, et
répétez en exagérant la syllabe forte. Tapez du doigt sur la table à chaque accent !}
\sonpaire{Ja\acc{pan}}{dʒəˈpæn}{le Japon}{Japa\acc{nese}}{ˌdʒæpəˈniːz}{japonais}
\sonpaire{\acc{Chi}na}{ˈtʃaɪnə}{la Chine}{Chi\acc{nese}}{tʃaɪˈniːz}{chinois}
\sonpaire{\acc{It}aly}{ˈɪtəli}{l'Italie}{I\acc{tal}ian}{ɪˈtæljən}{italien}
\sonpaire{\acc{Can}ada}{ˈkænədə}{le Canada}{Ca\acc{na}dian}{kəˈneɪdiən}{canadien}
\sonpaire{\acc{Ger}many}{ˈdʒɝːməni}{l'Allemagne}{\acc{Ger}man}{ˈdʒɝːmən}{allemand}
\sonpaire{\acc{Mex}ico}{ˈmeksɪkoʊ}{le Mexique}{\acc{Mex}ican}{ˈmeksɪkən}{mexicain}
\sonpaire{A\acc{mer}ica}{əˈmerɪkə}{l'Amérique}{A\acc{mer}ican}{əˈmerɪkən}{américain}
\sonpaire{\acc{Eng}land}{ˈɪŋɡlənd}{l'Angleterre}{\acc{Eng}lish}{ˈɪŋɡlɪʃ}{anglais}
\end{paires}

\begin{astuce}[Trois réflexes]
\begin{etapes}
\item Apprenez chaque mot nouveau \textbf{avec son accent} : regardez où se place le signe \phon{ˈ}.
\item N'accentuez pas la dernière syllabe « à la française » : on dit \en{\textbf{Ca}nada}, pas
  « Cana\textbf{da} ».
\item Les terminaisons \en{-ese} (\en{Japa\textbf{nese}}) portent l'accent ; avec \en{-ian},
  l'accent tombe juste avant (\en{Ca\textbf{na}dian}, \en{Bra\textbf{zil}ian}).
\end{etapes}
\end{astuce}

%% ---------------------------------------------------------------------
\rubrique{culture}{Le monde anglophone}

\begin{culture}[Une langue, des centaines de millions de locuteurs]
Près de 400 millions de personnes ont l'anglais pour langue maternelle, et plus d'un milliard
l'utilisent comme seconde langue : c'est la langue la plus parlée au monde. Elle est langue
officielle dans une soixantaine de pays, des États-Unis à l'Inde, en passant par le Canada,
l'Australie, l'Irlande, la Nouvelle-Zélande ou l'Afrique du Sud.
\end{culture}

\begin{culture}[Royaume-Uni, Grande-Bretagne ou Angleterre ?]
\begin{center}
\begin{tikzpicture}
  \fill[otTeal!12,rounded corners=2mm] (0,0) rectangle (150mm,40mm);
  \node[anchor=north west,font=\otSemi\fontsize{9}{11}\selectfont,text=otTeal] at (3mm,38mm)
    {\enb{the United Kingdom (the UK)} · le Royaume-Uni};
  \fill[otTeal!28,rounded corners=1.6mm] (3mm,3mm) rectangle (106mm,30mm);
  \node[anchor=north west,font=\otSemi\fontsize{8.4}{10}\selectfont,text=otTeal!80!black] at (6mm,28.5mm)
    {\enb{Great Britain} · la Grande-Bretagne};
  \foreach \n/\t/\cap/\x in {England/Angleterre/London/6,Scotland/Écosse/Edinburgh/40,Wales/pays de Galles/Cardiff/74}{%
    \fill[white,rounded corners=1.2mm] (\x mm,6mm) rectangle ++(30mm,15mm);
    \node[font=\otSemi\fontsize{9}{10}\selectfont] at (\x mm+15mm,16.5mm) {\enb{\n}};
    \node[font=\otRegular\fontsize{7.4}{9}\selectfont,text=otMuted] at (\x mm+15mm,12mm) {\t};
    \node[font=\otRegular\fontsize{7.4}{9}\selectfont,text=otTeal] at (\x mm+15mm,8.4mm) {\enb{\cap}};}
  \fill[white,rounded corners=1.2mm] (112mm,6mm) rectangle ++(34mm,15mm);
  \node[font=\otSemi\fontsize{9}{10}\selectfont] at (129mm,16.5mm) {\enb{Northern Ireland}};
  \node[font=\otRegular\fontsize{7.4}{9}\selectfont,text=otMuted] at (129mm,12mm) {Irlande du Nord};
  \node[font=\otRegular\fontsize{7.4}{9}\selectfont,text=otTeal] at (129mm,8.4mm) {\enb{Belfast}};
\end{tikzpicture}
\end{center}
Un Écossais ou un Gallois est \en{British}, mais il n'est pas \en{English} ! Pour ne vexer
personne, dites \en{British} quand vous parlez d'une personne du Royaume-Uni.
\end{culture}

\begin{paire}
\begin{culture}[Les États-Unis]
Cinquante États (\en{states}) et une capitale, Washington. New York est la plus grande ville du
pays, mais ce n'est pas la capitale ! On dit \en{the US}, \en{the USA} ou, à l'oral,
\en{the States}.
\end{culture}
\begin{culture}[\en{American}… ou pas ?]
Pour un habitant des États-Unis, on dit \en{American}. Mais attention : un Canadien ou un
Mexicain vit aussi en Amérique, et il n'aime pas qu'on le prenne pour un \en{American} !
\end{culture}
\end{paire}

%% ---------------------------------------------------------------------
\rubrique{pratique}{Exercices}

\exercice[ecrire,niveau=1]{Pays et nationalités}{Complétez avec la nationalité ou le pays qui manque.}
\begin{questions}[2]
\item Spain → \trou[26mm]{Spanish}
\item China → \trou[26mm]{Chinese}
\item Italy → \trou[26mm]{Italian}
\item Brazil → \trou[26mm]{Brazilian}
\item \trou[26mm]{Mexico} ← Mexican
\item \trou[26mm]{Ireland} ← Irish
\item \trou[26mm]{Canada} ← Canadian
\item \trou[26mm]{Germany} ← German
\end{questions}

\exercice[ecrire,niveau=1]{À la forme négative}{Récrivez chaque phrase à la forme négative, avec une
forme contractée.}
\begin{questions}[2]
\item I'm Italian. → \trou[40mm]{I'm not Italian.}
\item Jack is British. → \trou[40mm]{Jack isn't British.}
\item We're from Paris. → \trou[40mm]{We aren't from Paris.}
\item Emma is American. → \trou[40mm]{Emma isn't American.}
\item They're Japanese. → \trou[40mm]{They aren't Japanese.}
\item It's a big city. → \trou[40mm]{It isn't a big city.}
\end{questions}

\exercice[ecrire,niveau=2]{Posez la question}{Transformez chaque phrase en question.}
\begin{questions}
\item Léa is French. \hfill → \trou[70mm]{Is Léa French?}
\item You are from Brooklyn. \hfill → \trou[70mm]{Are you from Brooklyn?}
\item Emma is a beginner. \hfill → \trou[70mm]{Is Emma a beginner?}
\item They are Canadian. \hfill → \trou[70mm]{Are they Canadian?}
\item Mrs.\ Lopez is Mexican. \hfill → \trou[70mm]{Is Mrs. Lopez Mexican?}
\item Lyon is a big city. \hfill → \trou[70mm]{Is Lyon a big city?}
\end{questions}

\exercice[ecrire,niveau=1]{Réponses courtes}{Répondez par une réponse courte, d'après les dialogues.}
\begin{questions}[2]
\item Is Léa French? — \trou[32mm]{Yes, she is.}
\item Is Jack British? — \trou[32mm]{No, he isn't.}
\item Is Emma from London? — \trou[32mm]{Yes, she is.}
\item Are Jack and Emma from Paris? — \trou[32mm]{No, they aren't.}
\item Is Mrs.\ Lopez American? — \trou[32mm]{Yes, she is.}
\item Is Lyon a big city? — \trou[32mm]{Yes, it is.}
\end{questions}

\exercice[lire,ecrire,niveau=2]{Emma se présente}{Complétez avec \en{am}, \en{is}, \en{are},
\en{isn't} ou \en{aren't}.}
\begin{anglais}
\begin{tcolorbox}[otBase,colback=otNavyTint,colframe=otNavyTint,fontupper=\fontsize{11}{19}\selectfont]
Hello! My name \trou[14mm]{is} Emma Clarke. I \trou[14mm]{am} from London. I'm British, but my
best friend \trou[14mm]{isn't} British: she's Irish. My parents \trou[14mm]{are} in London.
They \trou[14mm]{aren't} in New York. And New York? It \trou[14mm]{is} a great city!
\par\vspace{1mm}{\fontsize{8.6}{10}\selectfont\color{otMuted}\foreignlanguage{french}{\textit{my best friend} : ma meilleure amie · \textit{great} : génial(e)}}
\end{tcolorbox}
\end{anglais}

\exercice[ecrire,niveau=2]{Attention aux majuscules}{Récrivez chaque phrase avec les majuscules
qui manquent.}
\begin{questions}
\item léa is french. → \trou[80mm]{Léa is French.}
\item jack speaks english and a little spanish. → \trou[80mm]{Jack speaks English and a little Spanish.}
\item emma is from london, in england. → \trou[80mm]{Emma is from London, in England.}
\item i'm not italian. → \trou[80mm]{I'm not Italian.}
\end{questions}

\exercice[ecouter,niveau=1,piste=2.10]{Où est l'accent ?}{Écoutez chaque mot, dit deux fois.
Soulignez la syllabe accentuée.}
\begin{questions}[4]
\item Ja · pan
\item Ja · pa · nese
\item It · a · ly
\item I · tal · ian
\item Can · a · da
\item Ca · na · di · an
\item A · mer · i · ca
\item Bra · zil
\end{questions}

\exercice[ecouter,ecrire,niveau=2,piste=2.11]{Qui vient d'où ?}{Quatre personnes se présentent.
Écoutez et complétez le tableau.}
\begin{center}\otTableau
\begin{tabular}{@{}l@{\hspace{12mm}}l@{\hspace{12mm}}l@{}}
\otTH{Prénom} & \otTH{Nationalité} & \otTH{Ville}\\\midrule
{\otSemi 1 · Ana}   & \enb{\trou[32mm]{Brazilian}} & \enb{\trou[36mm]{Rio de Janeiro}}\\[1mm]
{\otSemi 2 · Marco} & \enb{\trou[32mm]{Italian}}   & \enb{\trou[36mm]{Rome}}\\[1mm]
{\otSemi 3 · Yuki}  & \enb{\trou[32mm]{Japanese}}  & \enb{\trou[36mm]{Tokyo}}\\[1mm]
{\otSemi 4 · Liam}  & \enb{\trou[32mm]{Irish}}     & \enb{\trou[36mm]{Dublin}}\\\bottomrule
\end{tabular}\end{center}

\exercice[ecouter,ecrire,niveau=2,piste=2.12]{Dictée à trous}{Dans l'escalier, Léa croise Sam, son
voisin. Écoutez et complétez.}
\begin{anglais}
\noindent\begin{tabular}{@{}l@{\hspace{6mm}}l@{}}
\otNomDit{sam} & Hi, Léa! \trou{Are} you OK?\\[1.8mm]
\otNomDit{lea} & Yes, \trou{I'm} fine, thanks.\\[1.8mm]
\otNomDit{sam} & Are you \trou{from} Paris?\\[1.8mm]
\otNomDit{lea} & No, I'm \trou{not}. I'm from Lyon.\\[1.8mm]
\otNomDit{sam} & Is Lyon a \trou{big} city?\\[1.8mm]
\otNomDit{lea} & Yes, \trou{it} is. And you? Where are you from?\\[1.8mm]
\otNomDit{sam} & I'm \trou{from} here, from Brooklyn!\\[1.8mm]
\otNomDit{lea} & \trou{Really}? Great!\\
\end{tabular}
\end{anglais}

\exercice[ecrire,niveau=2]{Traduisez}{Traduisez en anglais. Utilisez les formes contractées.}
\begin{questions}[1][french]
\item Je viens de France.\lignes{1}
\item Tu es américain ? — Non. \textit{(réponse courte)}\lignes{1}
\item Elle n'est pas espagnole, elle est mexicaine.\lignes{1}
\item D'où viens-tu ?\lignes{1}
\item Je parle français et un peu anglais.\lignes{1}
\item Ils ne sont pas de Londres.\lignes{1}
\item Londres est une grande ville.\lignes{1}
\item Est-ce qu'Emma est britannique ? — Oui. \textit{(réponse courte)}\lignes{1}
\end{questions}

\exercice[lire,niveau=2]{Les nouveaux amis de Léa}{Lisez le texte de Léa, puis indiquez si chaque
phrase est vraie (V) ou fausse (F).}
\begin{anglais}
\begin{tcolorbox}[otBase,colback=white,colframe=otRule,boxrule=0.5pt,fontupper=\fontsize{10.8}{14.4}\selectfont]
Hello! I'm Léa. I'm French and I'm from Lyon. Now I'm in New York. Jack and Emma are my new
friends. Jack is American. He's from Brooklyn. Emma isn't American: she's British. She's from
London. Mrs.\ Lopez is American, but she speaks Spanish too. And me? I speak French, English
and a little Italian!
\par\vspace{1mm}{\fontsize{8.6}{10}\selectfont\color{otMuted}\foreignlanguage{french}{\textit{now} : maintenant · \textit{new friends} : nouveaux amis · \textit{too} : aussi}}
\end{tcolorbox}
\begin{questions}[2]
\item Léa is from Paris. \hfill\trou[9mm]{F}
\item Jack is from Brooklyn. \hfill\trou[9mm]{V}
\item Emma is American. \hfill\trou[9mm]{F}
\item Mrs.\ Lopez speaks Spanish. \hfill\trou[9mm]{V}
\item Léa speaks a little Italian. \hfill\trou[9mm]{V}
\item Léa is in Lyon now. \hfill\trou[9mm]{F}
\end{questions}
\end{anglais}

\exercice[ecouter,parler,niveau=3,piste=2.13]{À vous de parler !}{Emma veut tout savoir sur vous.
Répondez-lui \textbf{à voix haute} pendant chaque pause, puis comparez avec les réponses modèles
du corrigé.}
\begin{center}\otTableau
\begin{tabular}{@{}>{\otSemi\fontsize{8.6}{10}\selectfont\color{otNavy}}l@{\hspace{3mm}}>{\enOn}p{70mm}<{\enOff}@{\hspace{6mm}}p{78mm}@{}}
 & \otTH{Emma dit\ldots} & \otTH{Vous\ldots}\\\midrule
1. & Hi! I'm Emma. Where are you from? & dites de quel pays vous venez.\\
2. & Oh, really? What city are you from? & donnez le nom de votre ville.\\
3. & So, are you American? & répondez par une réponse courte, puis donnez votre nationalité.\\
4. & Are you a beginner in English? & répondez par une réponse courte.\\
5. & Don't worry, your English is good! & remerciez-la.\\\bottomrule
\end{tabular}\end{center}

%% ---------------------------------------------------------------------
\rubrique{bilan}{Faites le point}

\begin{essentiel}
\begin{minipage}[t]{0.56\linewidth}
\otEssTitre{Les phrases clés}
\renewcommand{\arraystretch}{1.32}
\begin{tabular}{@{}>{\enOn\otSemi\fontsize{9.4}{11}\selectfont}l<{\enOff}@{\hspace{3mm}}>{\fontsize{9.2}{11}\selectfont\color{white!78!otNavy}}l@{}}
Where are you from?            & D'où viens-tu ?\\
I'm from France.               & Je viens de France.\\
I'm French.                    & Je suis français(e).\\
Are you American?              & Tu es américain(e) ?\\
Yes, I am. / No, I'm not.      & Oui. / Non.\\
I speak French and English.    & Je parle français et anglais.\\
a little                       & un peu\\
That's right. / Really?        & C'est exact. / Vraiment ?\\
\end{tabular}
\end{minipage}\hfill
\begin{minipage}[t]{0.38\linewidth}
\otEssTitre{Négation}
{\enb{\fontsize{9.4}{12}\selectfont I'm \textbf{not} · he \textbf{isn't} · they \textbf{aren't}}}
\otEssTitre{Question}
{\enb{\fontsize{9.4}{12}\selectfont You are… → \textbf{Are you}…?\par Is she…? — Yes, she is.}}
\otEssTitre{Majuscules}
{\fontsize{9.2}{11.6}\selectfont Nationalités et langues : \en{French}, \en{English}, \en{Japanese}.\par}
\otEssTitre{Prononciation}
{\fontsize{9.2}{11.6}\selectfont Une syllabe forte par mot : \en{Ja\textbf{pan}}, \en{Japa\textbf{nese}}.\par}
\end{minipage}
\end{essentiel}

\begin{jesais}
\item dire de quel pays et de quelle ville je viens, et donner ma nationalité ;
\item demander à quelqu'un d'où il vient : \en{Where are you from?} ;
\item utiliser la forme négative de \en{be} : \en{I'm not}, \en{isn't}, \en{aren't} ;
\item poser une question avec \en{be} et répondre par une réponse courte ;
\item écrire les nationalités et les langues avec une majuscule ;
\item placer l'accent sur la bonne syllabe : \en{Japan}, \en{Japanese}.
\end{jesais}

\clearpage
\begin{testexpress}
\qcm[a]{« Je viens de France » se dit :}{\en{I'm from France.}}{\en{I come France.}}{\en{I'm of France.}}
\qcm[b]{La nationalité « japonais » se dit :}{\en{Japan}}{\en{Japanese}}{\en{Japanish}}
\qcm[a]{Complétez : \en{Jack \rule[-0.6mm]{9mm}{0.5pt} British.}}{\en{isn't}}{\en{aren't}}{\en{not}}
\qcm[b]{Quelle question est correcte ?}{\en{You are French?}}{\en{Are you French?}}{\en{Is you French?}}
\qcm[b]{Quelle réponse courte est correcte ?}{\en{Yes, I'm.}}{\en{Yes, I am.}}{\en{Yes, am I.}}
\qcm[c]{« Ils ne sont pas » se dit :}{\en{they isn't}}{\en{they not are}}{\en{they aren't}}
\qcm[b]{Quelle phrase est bien écrite ?}{\en{I speak english.}}{\en{I speak English.}}{\en{i speak English.}}
\qcm[c]{Dans \en{Japanese}, la syllabe accentuée est :}{\en{ja}}{\en{pa}}{\en{nese}}
\qcm[a]{« Le Royaume-Uni » se dit :}{\en{the UK}}{\en{UK}}{\en{the England}}
\qcm[a]{\en{Where's she from?} — }{\en{She's from Rome.}}{\en{She's Rome.}}{\en{She from Rome.}}
\end{testexpress}

\noindent{\fontsize{9.6}{12}\selectfont\color{otMuted}Niveau visé (CECRL, niveau A1) : « peut
répondre à des questions simples et en poser sur son origine et sa nationalité ». Les corrigés de
l'unité sont page~\pageref{cor-u2-1} ; les textes des pistes d'exercices, page~\pageref{trans-2.10}.\par}
